\chapter{Teilbarkeit, Division und Primzahlen}
\label{b10guide:ch:divisibility}

Alle Zahlen dieses Kapitels liegen in $\mathbb N$; zusätzliche
Positivitätsbedingungen stehen bei der jeweiligen Aussage.
Die folgenden Standardfolgen ergänzen die Arithmetik des Hauptbands.
Sie benötigen keine Erweiterung des Zahlbereichs und keine
Theorie von Restklassen. Vergleichsquellen sind
\cite{natgcd,natprime,natfactors,natinfinite}.

\section{Teilbarkeit und Division mit Rest}

\begin{definition}[Teilbarkeit]
\label{b10guide:def:divisibility}
Für $a,b\in\mathbb N$ bedeutet
\[
 a\mid b\quad\Longleftrightarrow\quad
 \exists c\in\mathbb N\;(b=ac).
\]
Man nennt $a$ einen \emph{Teiler} von $b$ und $b$ ein
\emph{Vielfaches} von $a$. Das Zeichen $a\nmid b$ verneint diese Aussage.
\end{definition}

\begin{theorem}[Grundregeln der natürlichen Teilbarkeit]
Für $a,b,c,u,v\in\mathbb N$ gilt
\[
\begin{gathered}
 1\mid a,\qquad a\mid a,\qquad a\mid0,\qquad
 0\mid a\ \Longleftrightarrow\ a=0,\\
 a\mid b\ \land\ b\mid c\ \Longrightarrow\ a\mid c,\\
 a\mid b\ \land\ a\mid c\ \Longrightarrow\ a\mid(bu+cv),\\
 a\mid b\ \land\ b\mid a\ \Longleftrightarrow\ a=b,\\
 a\mid b\ \land\ b>0\ \Longrightarrow\ a\le b.
\end{gathered}
\]
Ist $c>0$, so gilt außerdem $ca\mid cb\Longleftrightarrow a\mid b$.
\end{theorem}

\begin{theorem}[Division mit Rest]
\label{b10guide:thm:division}
Zu $a\in\mathbb N$ und $b\in\mathbb N_+$ gibt es genau ein Paar
$(q,r)\in\mathbb N\times\mathbb N$ mit
\[
 a=bq+r,\qquad 0\le r<b.
\]
\end{theorem}

\begin{definition}[Natürlicher Quotient und Rest]
Für $a\in\mathbb N$ und $b>0$ heißen die Zahlen aus
Satz~\ref{b10guide:thm:division}
\[
 a\mathbin{\mathrm{div}}b:=q,\qquad
 a\mathbin{\mathrm{mod}}b:=r.
\]
Division und Rest bei Divisor $0$ werden hier nicht definiert.
\end{definition}

\begin{theorem}[Charakterisierungen von Quotient und Rest]
Für $a,q,r\in\mathbb N$ und $b>0$ gilt
\[
\begin{aligned}
 q=a\mathbin{\mathrm{div}}b
 &\ \Longleftrightarrow\ bq\le a<b(q+1),\\
 r=a\mathbin{\mathrm{mod}}b
 &\ \Longleftrightarrow\ r<b\ \land\
       \exists t\in\mathbb N\;(a=bt+r),\\
 b\mid a&\ \Longleftrightarrow\ a\mathbin{\mathrm{mod}}b=0.
\end{aligned}
\]
Insbesondere ist $a\mathbin{\mathrm{div}}b=0$ genau für $a<b$;
in diesem Fall ist $a\mathbin{\mathrm{mod}}b=a$.
\end{theorem}

\begin{definition}[Gerade und ungerade natürliche Zahlen]
Eine Zahl $n$ heißt \emph{gerade}, wenn $n=2k$ für ein
$k\in\mathbb N$, und \emph{ungerade}, wenn $n=2k+1$ für ein
$k\in\mathbb N$ gilt.
\end{definition}

\begin{theorem}[Geradheit und Ungeradheit]
Jede natürliche Zahl ist genau eines von beiden: gerade oder ungerade.
Im jeweiligen Fall ist $k$ eindeutig. Für $m,n\in\mathbb N$ gilt:
$m+n$ ist genau dann gerade, wenn beide Zahlen gerade oder beide
ungerade sind; $mn$ ist genau dann ungerade, wenn beide Zahlen
ungerade sind. Der Nachfolger wechselt die Geradheit.
\end{theorem}

\section{Größter gemeinsamer Teiler und kleinstes gemeinsames Vielfaches}

\begin{definition}[ggT, kgV und Teilerfremdheit]
\label{b10guide:def:gcd}
Für $(a,b)\ne(0,0)$ ist $\operatorname{ggT}(a,b)$ die größte positive
natürliche Zahl, die sowohl $a$ als auch $b$ teilt.
Man setzt $\operatorname{ggT}(0,0)=0$.
Für $a,b>0$ ist $\operatorname{kgV}(a,b)$ das kleinste positive
gemeinsame Vielfache; falls $a=0$ oder $b=0$, setzt man
$\operatorname{kgV}(a,b)=0$.
Die Zahlen $a,b$ heißen \emph{teilerfremd}, wenn
$\operatorname{ggT}(a,b)=1$ gilt.
\end{definition}

\begin{theorem}[Existenz und arithmetische Charakterisierungen]
Die in Definition~\ref{b10guide:def:gcd} beschriebenen Zahlen
existieren und sind eindeutig. Für alle $a,b,d,m\in\mathbb N$ gilt
\[
\begin{aligned}
 d\mid\operatorname{ggT}(a,b)
 &\ \Longleftrightarrow\ d\mid a\ \land\ d\mid b,\\
 \operatorname{kgV}(a,b)\mid m
 &\ \Longleftrightarrow\ a\mid m\ \land\ b\mid m.
\end{aligned}
\]
\end{theorem}

\begin{theorem}[Randwerte und Produktformel]
Für $a,b\in\mathbb N$ gilt
\[
\begin{gathered}
 \operatorname{ggT}(a,0)=a,\qquad
 \operatorname{ggT}(a,1)=1,\qquad
 \operatorname{ggT}(a,b)=\operatorname{ggT}(b,a),\\
 \operatorname{kgV}(a,0)=0,\qquad
 \operatorname{kgV}(a,1)=a,\qquad
 \operatorname{kgV}(a,b)=\operatorname{kgV}(b,a),\\
 \operatorname{ggT}(a,b)\operatorname{kgV}(a,b)=ab.
\end{gathered}
\]
Insbesondere ist für teilerfremde $a,b$ das kgV gleich $ab$.
\end{theorem}

\begin{theorem}[Euklidischer Algorithmus]
\label{b10guide:thm:euclidean}
Für $a\in\mathbb N$ und $b>0$ gilt
\[
 \operatorname{ggT}(a,b)
 =\operatorname{ggT}\bigl(b,a\mathbin{\mathrm{mod}}b\bigr).
\]
Die wiederholte Ersetzung des Paars $(a,b)$ durch
$(b,a\mathbin{\mathrm{mod}}b)$, solange die zweite Zahl positiv ist,
endet nach endlich vielen Schritten bei einem Paar $(g,0)$.
Dessen erste Zahl ist der ggT des Ausgangspaars.
\end{theorem}

\begin{theorem}[Teilerfremdes Kürzen]
Für $a,b,c\in\mathbb N$ gilt
\[
 \operatorname{ggT}(a,b)=1\ \land\ a\mid bc
 \quad\Longrightarrow\quad a\mid c.
\]
Sind $a,b$ teilerfremd und teilen beide $c$, so teilt $ab$ die Zahl $c$.
Sind $a,b>0$ und $g=\operatorname{ggT}(a,b)$, so gibt es eindeutig
$u,v>0$ mit $a=gu$, $b=gv$; diese $u,v$ sind teilerfremd.
\end{theorem}

\section{Primzahlen und eindeutige Faktorisierung}

\begin{definition}[Primzahl und zusammengesetzte Zahl]
\label{b10guide:def:prime}
Eine natürliche Zahl $p$ heißt \emph{Primzahl}, wenn $p\ge2$ und
\[
 \forall d\in\mathbb N\;(d\mid p\ \Longrightarrow\ d=1\ \lor\ d=p).
\]
Eine Zahl $n\ge2$, die keine Primzahl ist, heißt
\emph{zusammengesetzt}. Die Zahlen $0$ und $1$ gehören zu keiner
dieser beiden Klassen.
\end{definition}

\begin{theorem}[Primteiler und Kriterium für Zusammengesetztheit]
Jede Zahl $n\ge2$ besitzt einen Primteiler. Für $n\ge2$ sind äquivalent:
\begin{enumerate}
 \item $n$ ist zusammengesetzt;
 \item es gibt $a,b\in\mathbb N$ mit $a,b\ge2$ und $n=ab$;
 \item es gibt einen Primteiler $p$ von $n$ mit $p^2\le n$.
\end{enumerate}
Insbesondere ist $n\ge2$ genau dann prim, wenn keine natürliche
Zahl $d$ mit $2\le d$ und $d^2\le n$ die Zahl $n$ teilt.
\end{theorem}

\begin{theorem}[Euklids Lemma]
\label{b10guide:thm:euclid-lemma}
Für jede Primzahl $p$ und $a,b\in\mathbb N$ gilt
\[
 p\mid ab\quad\Longleftrightarrow\quad p\mid a\ \lor\ p\mid b.
\]
Für $p\ge2$ charakterisiert diese Eigenschaft, für alle $a,b$
gefordert, die Primzahlen. Ist $k>0$, so gilt für eine Primzahl $p$
auch $p\mid a^k\Longleftrightarrow p\mid a$.
\end{theorem}

\begin{theorem}[Hauptsatz der elementaren Arithmetik]
\label{b10guide:thm:factorization}
Jede natürliche Zahl $n>1$ besitzt genau eine Darstellung
\[
 n=p_1p_2\cdots p_k,\qquad
 k\ge1,\qquad p_1\le p_2\le\cdots\le p_k,
\]
in der alle $p_i$ Primzahlen sind. Die Eindeutigkeit betrifft
sowohl $k$ als auch jeden einzelnen Faktor; Wiederholungen sind erlaubt.
Die Zahl $1$ wird durch das leere Produkt dargestellt.
Die Zahl $0$ besitzt keine solche Darstellung.
\end{theorem}

\begin{theorem}[Unbeschränktheit der Primzahlen]
\label{b10guide:thm:infinite-primes}
Zu jeder natürlichen Zahl $n$ gibt es eine Primzahl $p$ mit $p>n$.
Insbesondere gibt es keine größte Primzahl.
Die einzige gerade Primzahl ist $2$.
\end{theorem}

\section{Rechnen mit natürlichen Quotienten und Resten}

Die folgenden Regeln betreffen die natürlichen Zahlenwerte von
$\mathrm{div}$ und $\mathrm{mod}$ bei positiven Divisoren
\cite{natdivmod}.

\begin{theorem}[Verschiebung und geschachtelte Quotienten]
\label{b10guide:thm:quotient-rules}
Für $a,t\in\mathbb N$ und $b,c\in\mathbb N_+$ gilt
\[
 \begin{aligned}
 (a+bt)\mathbin{\mathrm{div}}b
   &=(a\mathbin{\mathrm{div}}b)+t,\\
 (a+bt)\mathbin{\mathrm{mod}}b&=a\mathbin{\mathrm{mod}}b,\\
 (a\mathbin{\mathrm{div}}b)\mathbin{\mathrm{div}}c
   &=a\mathbin{\mathrm{div}}(bc).
 \end{aligned}
\]
\end{theorem}

\begin{theorem}[Reste von Summen, Produkten und Potenzen]
\label{b10guide:thm:remainder-rules}
Für $a,b,k\in\mathbb N$ und $m\in\mathbb N_+$ gilt
\[
 \begin{aligned}
 (a\mathbin{\mathrm{mod}}m)\mathbin{\mathrm{mod}}m
   &=a\mathbin{\mathrm{mod}}m,\\
 (a+b)\mathbin{\mathrm{mod}}m
   &=\bigl((a\mathbin{\mathrm{mod}}m)
          +(b\mathbin{\mathrm{mod}}m)\bigr)\mathbin{\mathrm{mod}}m,\\
 (ab)\mathbin{\mathrm{mod}}m
   &=\bigl((a\mathbin{\mathrm{mod}}m)
          (b\mathbin{\mathrm{mod}}m)\bigr)\mathbin{\mathrm{mod}}m,\\
 a^k\mathbin{\mathrm{mod}}m
   &=\bigl((a\mathbin{\mathrm{mod}}m)^k\bigr)\mathbin{\mathrm{mod}}m.
 \end{aligned}
\]
Die Potenzregel umfaßt $k=0$ mit der Konvention $0^0=1$.
\end{theorem}

\section{Primexponenten und ihre Rechenregeln}

Die folgenden Aussagen beschreiben die Primfaktorzerlegung durch
natürliche Exponenten; sie führen keinen neuen Zahlbereich ein
\cite{natvaluation}.

\begin{definition}[Primexponent einer positiven natürlichen Zahl]
\label{b10guide:def:prime-exponent}
Für eine Primzahl $p$ und $n\in\mathbb N_+$ heißt die eindeutig
bestimmte natürliche Zahl $e$ mit
\[
 p^e\mid n,\qquad p^{e+1}\nmid n
\]
der \emph{Primexponent von $p$ in $n$}, geschrieben $v_p(n)$.
Ein Wert $v_p(0)$ wird hier nicht definiert.
\end{definition}

\begin{theorem}[Primexponenten von Produkten und Potenzen]
\label{b10guide:thm:prime-exponent-laws}
Für jede Primzahl $p$, positive natürliche Zahlen $a,b$ und
$e,k\in\mathbb N$ gilt
\[
 \begin{gathered}
 v_p(1)=0,\qquad v_p(p)=1,\\
 v_p(ab)=v_p(a)+v_p(b),\qquad v_p(a^k)=k\,v_p(a),\\
 p^e\mid a\quad\Longleftrightarrow\quad e\le v_p(a).
 \end{gathered}
\]
Insbesondere gilt $v_p(a)=0$ genau dann, wenn $p\nmid a$.
\end{theorem}

\begin{theorem}[Teilbarkeit und Gleichheit durch Primexponenten]
\label{b10guide:thm:prime-exponent-comparison}
Für $a,b\in\mathbb N_+$ gilt
\[
 \begin{aligned}
 a\mid b
 &\ \Longleftrightarrow\ \forall p\text{ prim}\;v_p(a)\le v_p(b),\\
 a=b
 &\ \Longleftrightarrow\ \forall p\text{ prim}\;v_p(a)=v_p(b),\\
 \operatorname{ggT}(a,b)=1
 &\ \Longleftrightarrow\ \forall p\text{ prim}\;
       \bigl(v_p(a)=0\ \lor\ v_p(b)=0\bigr).
 \end{aligned}
\]
\end{theorem}

\begin{theorem}[Erkennen vollständiger Potenzen]
\label{b10guide:thm:perfect-power-criterion}
Für $n,r\in\mathbb N_+$ gilt
\[
 \exists a\in\mathbb N\;(n=a^r)
 \quad\Longleftrightarrow\quad
 \forall p\text{ prim}\;(r\mid v_p(n)).
\]
Falls eine solche Zahl $a$ existiert, ist sie eindeutig und positiv.
Für $n=1$ sind alle Primexponenten $0$ und es gilt $a=1$.
\end{theorem}
\begin{theorem}[Primexponenten von ggT und kgV]
\label{b10guide:thm:prime-exponent-gcd-lcm}
Seien $a,b\in\mathbb N_+$, $p$ eine Primzahl und
$e=v_p(a)$, $f=v_p(b)$. Dann gilt
\[
 \Bigl(v_p\bigl(\operatorname{ggT}(a,b)\bigr),
       v_p\bigl(\operatorname{kgV}(a,b)\bigr)\Bigr)
 =\begin{cases}(e,f),&e\le f,\\(f,e),&f<e.\end{cases}
\]
\end{theorem}

\begin{theorem}[Primexponent einer Summe]
\label{b10guide:thm:prime-exponent-sum}
Seien $p$ eine Primzahl und $a,b\in\mathbb N_+$.
Setze $e=v_p(a)$ und $f=v_p(b)$.
Dann ist $v_p(a+b)$ mindestens so groß wie der kleinere der beiden
Werte $e,f$. Für $e\ne f$ gilt genauer
\[
 v_p(a+b)=\begin{cases}e,&e<f,\\f,&f<e.\end{cases}
\]
Bei $e=f$ gilt $v_p(a+b)\ge e$; eine Gleichheit wird in diesem Fall
nicht behauptet \cite{natvalsum}.
\end{theorem}

\section{Ausgleichsgleichungen und Potenzdifferenzen}

\begin{theorem}[Bézouts natürliche Ausgleichsgleichung]
\label{b10guide:thm:natural-bezout}
Für $a\in\mathbb N_+$ und $b\in\mathbb N$ gibt es
$x\in\mathbb N_+$ und $y\in\mathbb N$ mit
\[
 ax=by+\operatorname{ggT}(a,b).
\]
Insbesondere gilt unter denselben Voraussetzungen
\[
 \operatorname{ggT}(a,b)=1
 \quad\Longleftrightarrow\quad
 \exists x\in\mathbb N_+\;\exists y\in\mathbb N\;(ax=by+1).
\]
Auch $b=0$ ist zugelassen; dann erfüllt $x=1$, $y=0$ die erste
Gleichung \cite{natbezout}.
\end{theorem}

\begin{theorem}[ggT und Teilbarkeit von Potenzdifferenzen]
\label{b10guide:thm:gcd-power-differences}
Für $a\ge2$ und $m,n\in\mathbb N$ gilt
\[
 \operatorname{ggT}(a^m-1,a^n-1)
    =a^{\operatorname{ggT}(m,n)}-1.
\]
Außerdem gilt
\[
 (a^m-1)\mid(a^n-1)\quad\Longleftrightarrow\quad m\mid n.
\]
Alle Differenzen sind natürliche Differenzen. Die Aussagen schließen
$m=0$ und $n=0$ mit $a^0-1=0$ und
$\operatorname{ggT}(0,0)=0$ ein \cite{natgcd}.
\end{theorem}


\begin{theorem}[Primexponenten von Potenzdifferenzen -- LTE]
\label{b10guide:thm:lte}
Seien $a,b,n\in\mathbb N$ mit $a>b\ge1$ und $n\ge1$.
Für eine ungerade Primzahl $p$ mit $p\mid(a-b)$ und $p\nmid a$ gilt
\[
 v_p(a^n-b^n)=v_p(a-b)+v_p(n).
\]
Sind $a$ und $b$ ungerade und ist $n$ gerade, so gilt
\[
 v_2(a^n-b^n)+1=v_2(a-b)+v_2(a+b)+v_2(n).
\]
Die Bedingungen $a>b$ und $n\ge1$ gewährleisten, daß sämtliche
Argumente der Primexponenten positiv sind. Die zweite Formel
verwendet nur natürliche Summen \cite{natlte}.
\end{theorem}
